% Part: turing-machines
% Chapter: machines-computations
% Section: turing-machines

\documentclass[../../../include/open-logic-section]{subfiles}

\begin{document}

\olfileid{tur}{mac}{tur}
\olsection{টুরিং যন্ত্র}

\begin{explain}
কোনটি টুরিং যন্ত্র হিসেবে গণ্য হবে তার আনুষ্ঠানিক সংজ্ঞাটি দেখতে বিমূর্ত,
কিন্তু আসলে সরল: টুরিং যন্ত্রের কাজকর্ম নির্দিষ্ট করতে যত তথ্য দরকার,
সেগুলি শুধু একটি গাণিতিক গঠনে একত্র করা হয়। এর মধ্যে আছে (১) যন্ত্রটি
কোন কোন দশায় থাকতে পারে, (২) ফিতায় কোন প্রতীকগুলি থাকতে পারে, (৩)
যন্ত্রটি কোন দশায় শুরু করবে এবং (৪) যন্ত্রটির নির্দেশ-সেট কী।
\end{explain}

\begin{defn}[টুরিং যন্ত্র]
একটি \emph{টুরিং যন্ত্র} $M$ হলো একটি চতুষ্টয় $\langle Q, \Sigma, q_0,
\delta\rangle$, যার উপাদানগুলি হলো
\begin{enumerate}
\item \emph{দশা}-গুলির একটি সসীম সেট~$Q$,
\item একটি সসীম \emph{বর্ণমালা}~$\Sigma$, যার মধ্যে $\TMendtape$ ও
  $\TMblank$ আছে,
\item একটি \emph{আরম্ভিক দশা}~$q_0 \in Q$,
\item একটি সসীম \emph{নির্দেশ-সেট}~$\delta\colon Q \times \Sigma
  \pto Q \times \Sigma \times \{\TMleft, \TMright, \TMstay\}$।
\end{enumerate}
আংশিক অপেক্ষক~$\delta$-কে $M$-এর \emph{অবস্থান্তর অপেক্ষক}-ও বলা হয়।
\end{defn}

\begin{explain}
আমরা ধরে নিই, ফিতাটি শুধু এক দিকেই অসীম। তাই ফিতার বাঁ প্রান্তের চিহ্ন
হিসেবে একটি বিশেষ প্রতীক~$\TMendtape$ নির্দিষ্ট করা সুবিধাজনক। এতে
টুরিং-যন্ত্রের প্রোগ্রাম বুঝতে পারে, কখন ফিতার বাইরে চলে যাওয়ার “বিপদে”
পড়েছে। আমরা ধরে নিতে পারতাম যে এই প্রতীকটির উপর কখনও অন্য কিছু লেখা
হয় না; অর্থাৎ $\delta(q,\TMendtape)$ সংজ্ঞায়িত হলে
$\delta(q,\TMendtape) = \tuple{q', \TMendtape, x}$। কোনো কোনো পাঠ্যবইয়ে
এই অনুমান করা হয়; আমরা করি না। টুরিং যন্ত্র নির্মাণের সময় সতর্ক থেকে
শুধু নিশ্চিত করতে পারো যে সেটি কখনও~$\TMendtape$-এর উপর অন্য কিছু লেখে
না। তাছাড়া এমন ক্ষেত্র আছে, যেখানে তার উপর লেখার অনুমতি কিছু সুবিধাজনক
নমনীয়তা দেয়।
\end{explain}

\begin{ex}
\emph{জোড়-যন্ত্র:} জোড়-যন্ত্রটি আনুষ্ঠানিকভাবে
$\tuple{Q, \Sigma, q_0, \delta}$ চতুষ্টয়, যেখানে
\begin{align*}
Q & = \{ q_0, q_1 \} \\
\Sigma & = \{ \TMendtape, \TMblank, \TMstroke \}, \\
\delta(q_0, \TMstroke) & = \tuple{q_1, \TMstroke, \TMright},\\
\delta(q_1, \TMstroke) & = \tuple{q_0, \TMstroke, \TMright},\\
\delta(q_1, \TMblank)  & = \tuple{q_1, \TMblank, \TMright}.
\end{align*}
\end{ex}

\end{document}
